
\documentclass[11pt,a4paper]{article}

\usepackage[T1]{fontenc}
\usepackage[utf8]{inputenc}
\usepackage[polish]{babel}
\usepackage{lmodern}
\usepackage{geometry}
\usepackage{booktabs}
\usepackage{longtable}
\usepackage{array}
\usepackage{xcolor}
\usepackage{enumitem}
\usepackage{hyperref}
\usepackage{ragged2e}
\usepackage{caption}
\usepackage{fancyhdr}

\geometry{margin=2.2cm, bottom=2.8cm}
\hypersetup{colorlinks=true, linkcolor=black, urlcolor=blue!50!black}

\newcolumntype{L}[1]{>{\RaggedRight\arraybackslash}p{#1}}
\newcommand{\kryt}{\textcolor{red!70!black}{\textbf{[K]}}}
\newcommand{\testid}[1]{\texttt{#1}}

\pagestyle{fancy}
\fancyhf{}
\fancyhead[L]{\small bApka --- katalog testów jednostkowych (43)}
\fancyhead[R]{\small logika biznesowa i reguły dostępu}
\fancyfoot[C]{\thepage}

\title{Katalog przypadków testowych\\[0.4em]
\large Testy jednostkowe logiki biznesowej i reguł dostępu\\[0.3em]
\normalsize Plan rdzenny: 43 przypadków}
\author{Projekt bApka (wersja robocza pod rozdział pracy inżynierskiej)}
\date{\today}

\begin{document}
\maketitle

\begin{abstract}
\noindent
Dokument zawiera \textbf{43 rdzenne} przypadki testów jednostkowych
(wcześniejszy katalog roboczy został zawężony; z E usunięto FCM dedup i SOS create).
Priorytet: bezpieczeństwo dostępu (IDOR), integralność dawek/magazynu,
auth (credentials / rate limit), SOS oraz krytyczne ścieżki UI seniora.
Oznaczenie \kryt{} = przypadek krytyczny.
Warstwa: Jest + mocki --- bez kontraktu REST, live PostgreSQL i E2E sprzętowego.
\end{abstract}

\tableofcontents
\newpage

\section{Konwencje i zakres}

\begin{itemize}[leftmargin=*]
  \item \textbf{Liczba planowana:} 43 ID (A--H), rozłożone równomiernie z naciskiem
        na przypadki krytyczne.
  \item \textbf{Usunięte z planu:} m.in.\ trywialne NotFound, duplikaty CRUD,
        drobne normalizacje, drugorzędne utils FE, stepper/HugeButton, IoT stub.
  \item \textbf{Poza zakresem:} testy kontraktowe REST, integracja DB, E2E, obciążenie.
\end{itemize}

\subsection*{Rozkład 43 przypadków}
\begin{center}
\begin{tabular}{@{}lc@{}}
\toprule
Kategoria & Liczba ID \\
\midrule
A Reguły dostępu / IDOR & 7 \\
B Auth / credentials & 5 \\
C Harmonogramy / dawki & 8 \\
D Magazyn & 5 \\
E Pairing / profil & 3 \\
F Powiadomienia / cron & 5 \\
G Frontend (bramki / routing) & 5 \\
H Interakcje seniora & 5 \\
\midrule
\textbf{Razem} & \textbf{43} \\
\bottomrule
\end{tabular}
\end{center}

\newcommand{\casetablehead}{%
  \toprule
  ID & Skrót & Funkcja / plik & Opis: co testujemy i po co & Oczekiwany wynik \\
  \midrule
}
\newcommand{\casetablefoot}{\bottomrule}

% =============================================================================
\section{A. Reguły dostępu (7)}
% =============================================================================

\begin{longtable}{@{}L{1.6cm}L{2.4cm}L{3.2cm}L{5.2cm}L{3.4cm}@{}}
\caption{Kategoria A --- 7 przypadków}\\
\casetablehead
\endfirsthead
\caption[]{Kategoria A (kontynuacja)}\\
\casetablehead
\endhead
\casetablefoot
\endfoot

\testid{A-01} &
Self-access &
\texttt{assertCanAccessUser} &
Właściciel danych ma dostęp do własnego \texttt{userId} bez odpytywania Connection. &
OK; \texttt{findFirst} nie wołane. \\
\midrule

\testid{A-02} \kryt &
Opiekun sparowany &
\texttt{assertCanAccessUser} &
Sparowany caretaker może zarządzać danymi podopiecznego. &
OK przy \texttt{isPaired=true}. \\
\midrule

\testid{A-03} \kryt &
Obcy / IDOR &
\texttt{assertCanAccessUser} &
Brak pairingu $\to$ odmowa (klasyczny IDOR na UUID). &
\texttt{ForbiddenException}. \\
\midrule

\testid{A-06} \kryt &
Harmonogram obcy &
\texttt{assertCanAccessSchedule} &
Schedule istnieje, aktor niesparowany $\to$ 403. &
Forbidden po odczycie właściciela. \\
\midrule

\testid{A-08} \kryt &
Magazyn obcy &
\texttt{assertCanAccessInventory} &
Inventory istnieje, brak pairingu $\to$ 403. &
\texttt{ForbiddenException}. \\
\midrule

\testid{A-09} \kryt &
Settings seniora &
\texttt{updateDependentSettingsByCaretaker} &
Tylko sparowany caretaker + rola DEPENDENT. &
403 / NotFound gdy warunki niespełnione. \\
\midrule

\testid{A-10} \kryt &
Unpair / delete &
\texttt{unpairDependent}, \texttt{deleteDependentAccount} &
Operacje destrukcyjne tylko dla sparowanego opiekuna. &
Obcy $\to$ Forbidden. \\

\end{longtable}

\noindent\textit{Usunięte z A:} A-04 (HYBRID=self --- wariant A-01/A-03), A-05/A-07 (same NotFound),
A-11 (listy --- pokryte wzorcem access), A-12 (JWT --- warstwa auth).

% =============================================================================
\section{B. Auth (5)}
% =============================================================================

\begin{longtable}{@{}L{1.6cm}L{2.4cm}L{3.2cm}L{5.2cm}L{3.4cm}@{}}
\caption{Kategoria B --- 5 przypadków}\\
\casetablehead
\endfirsthead
\caption[]{Kategoria B (kontynuacja)}\\
\casetablehead
\endhead
\casetablefoot
\endfoot

\testid{B-02} \kryt &
Wrogie e-maile &
\texttt{isValidAuthEmail} &
Odrzucenie formatów wrogich / spoof / CRLF / granic długości. &
\texttt{false} / issue. \\
\midrule

\testid{B-03} &
Granice hasła &
\texttt{isValidAuthPassword} &
Min/max długość (bcrypt) + znaki kontrolne. &
Za krótkie/długie odrzucone. \\
\midrule

\testid{B-04} &
Duplikat e-mail &
\texttt{AuthService.register} &
Ponowna rejestracja tego samego adresu. &
\texttt{BadRequestException}. \\
\midrule

\testid{B-07} &
Login OK &
\texttt{login} / \texttt{validateUser} &
Poprawne credentials $\to$ JWT + skrót profilu. &
\texttt{access\_token} + user. \\
\midrule

\testid{B-09} \kryt &
Rate limit &
\texttt{AuthRateLimitGuard} &
Anty brute-force: 10/min login, 5/min register. &
429 po przekroczeniu. \\

\end{longtable}

\noindent\textit{Usunięte z B:} B-01 (trywialna normalizacja), B-05 (default role),
B-06 (wariant B-07), B-08 (Google --- opcjonalny tor), B-10 (konfig env).

% =============================================================================
\section{C. Harmonogramy i dawki (8)}
% =============================================================================

\begin{longtable}{@{}L{1.6cm}L{2.4cm}L{3.2cm}L{5.2cm}L{3.4cm}@{}}
\caption{Kategoria C --- 8 przypadków}\\
\casetablehead
\endfirsthead
\caption[]{Kategoria C (kontynuacja)}\\
\casetablehead
\endhead
\casetablefoot
\endfoot

\testid{C-01} &
ONCE tylko w dniu &
\texttt{scheduleAppliesToDate} &
Terapia jednorazowa nie wycieka na inne dni. &
Tylko \texttt{startDate}. \\
\midrule

\testid{C-04} \kryt &
Spójność BE$\leftrightarrow$FE &
oba \texttt{scheduleAppliesToDate} &
Ten sam zestaw dat $\to$ identyczny wynik BE/FE. &
Parametrycznie równe. \\
\midrule

\testid{C-05} \kryt &
TAKEN w oknie &
\texttt{markDose} &
Potwierdzenie w $\pm$5\,min $\to$ TAKEN. &
Status TAKEN + \texttt{APP\_SENIOR}. \\
\midrule

\testid{C-06} \kryt &
LATE po oknie &
\texttt{markDose} &
Po oknie klik ``wziąłem'' $\to$ LATE. &
Status LATE. \\
\midrule

\testid{C-07} &
MISSED &
\texttt{markDose(..., MISSED)} &
Jawne pominięcie; bez zużycia tabletek. &
MISSED; brak dekrementacji. \\
\midrule

\testid{C-10} \kryt &
Europe/Warsaw &
\texttt{getLocalDayRange} &
Granice dnia lokalnego, nie UTC. &
Poprawny zakres doby. \\
\midrule

\testid{C-14} \kryt &
Brak logu = MISSED &
\texttt{buildDoseStatsPayload} &
Po deadline slot bez logu nie jest ``OK''. &
MISSED w stats. \\
\midrule

\testid{C-15} &
Create + access &
\texttt{SchedulesService.create} &
Tworzenie harmonogramu wymaga dostępu do userId. &
403 lub OK z defaultami. \\

\end{longtable}

\noindent\textit{Usunięte z C:} C-02/C-03 (pokryte zestawem C-04), C-08/C-09 (mechanika update/create),
C-11 (\texttt{normalizeHm}), C-12/C-13 (adherence UI), C-16 (cascade --- bliżej D).

% =============================================================================
\section{D. Magazyn (5)}
% =============================================================================

\begin{longtable}{@{}L{1.6cm}L{2.4cm}L{3.2cm}L{5.2cm}L{3.4cm}@{}}
\caption{Kategoria D --- 5 przypadków}\\
\casetablehead
\endfirsthead
\caption[]{Kategoria D (kontynuacja)}\\
\casetablehead
\endhead
\casetablefoot
\endfoot

\testid{D-01} &
Tylko MEDICATION &
\texttt{isMedicationInventory} &
Pomiary/ćwiczenia nie są zapasem tabletek. &
\texttt{true} tylko dla leku. \\
\midrule

\testid{D-02} &
Dni do wyczerpania &
\texttt{calculateInventoryDepletion} &
Szacunek \texttt{daysLeft} z recurring schedules. &
Poprawna arytmetyka / wysoki daysLeft. \\
\midrule

\testid{D-03} \kryt &
Pierwsze TAKEN odejmuje &
\texttt{markDose} + inventory &
Pierwsze PENDING$\to$TAKEN/LATE zmniejsza zapas. &
\texttt{inventory.update}. \\
\midrule

\testid{D-04} \kryt &
Idempotencja &
\texttt{markDose} &
Drugie potwierdzenie nie odejmuje ponownie. &
Brak drugiego update odejmującego. \\
\midrule

\testid{D-05} &
MISSED bez odejmowania &
\texttt{markDose} &
Pominięcie nie zużywa opakowania. &
Brak dekrementacji. \\

\end{longtable}

% =============================================================================
\section{E. Pairing i profil (3)}
% =============================================================================

\begin{longtable}{@{}L{1.6cm}L{2.4cm}L{3.2cm}L{5.2cm}L{3.4cm}@{}}
\caption{Kategoria E --- 3 przypadki}\\
\casetablehead
\endfirsthead
\caption[]{Kategoria E (kontynuacja)}\\
\casetablehead
\endhead
\casetablefoot
\endfoot

\testid{E-01} \kryt &
Pair tylko DEPENDENT &
\texttt{pairWithPin} &
HYBRID/CARETAKER nie paruje się jako dependent. &
Odrzucenie złej roli. \\
\midrule

\testid{E-02} \kryt &
Zły / zużyty PIN &
\texttt{pairWithPin} &
PIN sparowany lub nieistniejący nie łączy kont. &
Błąd biznesowy. \\
\midrule

\testid{E-03} &
Pair sukces &
\texttt{pairWithPin} &
Po sukcesie \texttt{isPaired=true}, PIN wyczyszczony. &
Connection sparowane. \\

\end{longtable}

\noindent\textit{Usunięte z E:} FCM dedup (\texttt{updateFcmToken}), SOS create self
(\texttt{createSosLog}) --- pokryte częściowo w F / istniejący \texttt{users.service.spec}.

% =============================================================================
\section{F. Powiadomienia i cron (5)}
% =============================================================================

\begin{longtable}{@{}L{1.6cm}L{2.4cm}L{3.2cm}L{5.2cm}L{3.4cm}@{}}
\caption{Kategoria F --- 5 przypadków}\\
\casetablehead
\endfirsthead
\caption[]{Kategoria F (kontynuacja)}\\
\casetablehead
\endhead
\casetablefoot
\endfoot

\testid{F-01} &
Treści powiadomień &
\texttt{activity-notification-messages} &
Copy zależny od kind / late / języka. &
Macierz asercji. \\
\midrule

\testid{F-02} \kryt &
SOS do caretakerów &
\texttt{notifySos} &
Alarm tylko do sparowanych tokenów. &
Push do paired; cleanup stale. \\
\midrule

\testid{F-03} \kryt &
Progi magazynu &
\texttt{processInventoryAlerts} &
Low/empty + flagi ``już wysłano'' (bez spamu). &
Pierwszy raz wyślij; drugi skip. \\
\midrule

\testid{F-04} &
Reminder minuty &
\texttt{handleMedicationReminders} &
Slot = minuta zegara, brak logu $\to$ PENDING + push. &
Create + reminder. \\
\midrule

\testid{F-05} \kryt &
PENDING$\to$MISSED &
\texttt{markMissedDoses} &
Po oknie 5\,min PENDING staje się MISSED + notify. &
Update + powiadomienia. \\

\end{longtable}

% =============================================================================
\section{G. Frontend bramki (5)}
% =============================================================================

\begin{longtable}{@{}L{1.6cm}L{2.4cm}L{3.2cm}L{5.2cm}L{3.4cm}@{}}
\caption{Kategoria G --- 5 przypadków}\\
\casetablehead
\endfirsthead
\caption[]{Kategoria G (kontynuacja)}\\
\casetablehead
\endhead
\casetablefoot
\endfoot

\testid{G-01} \kryt &
Assert dependent &
\texttt{assertCaretakerDependent} &
UI nie buduje żądań na obcym UUID. &
Błąd gdy unlinked. \\
\midrule

\testid{G-02} \kryt &
DEPENDENT$\to$PIN &
\texttt{postAuthRouting} &
Niesparowany dependent idzie na \texttt{/enter-pin}. &
Właściwy route. \\
\midrule

\testid{G-03} &
Dose stats FE &
\texttt{doseStats.ts} &
Spójność adherence z backend helpers. &
Te same statusy/\%. \\
\midrule

\testid{G-04} &
Push router &
\texttt{incomingPushRouter} &
Decyzje sync/nawigacji po typie push. &
Właściwy emit/route. \\
\midrule

\testid{G-05} &
Mood schedule &
\texttt{computeMoodScheduleState} &
Kafelek nastroju: disabled / active / inactive. &
Poprawny kind. \\

\end{longtable}

% =============================================================================
\section{H. Interakcje seniora (5)}
% =============================================================================

\begin{longtable}{@{}L{1.6cm}L{2.4cm}L{3.2cm}L{5.2cm}L{3.4cm}@{}}
\caption{Kategoria H --- 5 przypadków}\\
\casetablehead
\endfirsthead
\caption[]{Kategoria H (kontynuacja)}\\
\casetablehead
\endhead
\casetablefoot
\endfoot

\testid{H-01} \kryt &
Klik dawki due &
dashboard + \texttt{markTaken} &
W oknie $\to$ modal TAK $\to$ API markTaken. &
API raz + completion. \\
\midrule

\testid{H-02} \kryt &
Klik upcoming &
dashboard &
Poza oknem nie wolno markTaken. &
Brak wywołania API. \\
\midrule

\testid{H-03} &
Mood pick &
\texttt{MoodPickerModal} / gate &
Wybór buźki $\to$ \texttt{updateMood}. &
Właściwy mood. \\
\midrule

\testid{H-04} \kryt &
SOS confirm &
\texttt{SeniorConfirmModal} &
SOS dopiero po potwierdzeniu. &
Anuluj bez API; TAK raz. \\
\midrule

\testid{H-05} &
Modal TAK/NIE &
\texttt{SeniorConfirmModal} &
Kontrakt callbacków confirm/cancel. &
Właściwy callback. \\

\end{longtable}

% =============================================================================
\section{Zestawienie krytycznych i stan implementacji}
% =============================================================================

\begin{table}[h]
\centering
\caption{Przypadki \kryt{} w planie 43}
\begin{tabular}{@{}ll@{}}
\toprule
ID & Temat \\
\midrule
A-02, A-03, A-06, A-08, A-09, A-10 & IDOR / pairing \\
B-02, B-09 & credentials / rate limit \\
C-04, C-05, C-06, C-10, C-14 & dawki / czas / stats \\
D-03, D-04 & magazyn idempotentny \\
E-01, E-02 & PIN pairing \\
F-02, F-03, F-05 & SOS / alerty / MISSED \\
G-01, G-02 & FE gate / PIN route \\
H-01, H-02, H-04 & klik seniora / SOS \\
\bottomrule
\end{tabular}
\end{table}

\begin{itemize}[leftmargin=*]
  \item Zaimplementowane pliki kategorii: \texttt{category-a-access-rules.spec.ts},
        \texttt{category-b-auth-identity.spec.ts}, \texttt{category-c-schedules-doses.spec.ts},
        \texttt{category-d-inventory.spec.ts}, \texttt{category-e-pairing-profile.spec.ts},
        \texttt{category-f-notifications-cron.spec.ts},
        \texttt{category-g-frontend-gates.test.ts},
        \texttt{category-h-senior-interactions.test.ts} (frontend)
        (mogą zawierać \emph{więcej} niż ID z planu 43 --- nadmiar to pokrycie dodatkowe).
  \item Do zrobienia wg planu: --- (A--H wg katalogu 43).
\end{itemize}

\end{document}
